\title{Диагностика заболеваний дынь по визуальным признакам,\\ выращиваемых в условиях нестабильного содержания\\ питательных веществ в почве}

\taskdate{01.01.2025}

\projecttasks{%
  \projecttask{\bfseries\projecttasknum}{\bfseries Стандартная задача 1}{}{}{}%
  \projecttask{\projectsubtasknum}%
  {В этом разделе будет представлен стандартный текст-заполнитель для задач.}%
  {Результаты стандартной задачи}%
  {}{\signat[4pt]{.12}{StandardUser}{01.01.2025}}%
}

\taskliterature{%
\nocite{StandardRef1}
\nocite{StandardRef2}
}

% # Подписи

% Для простановки подписи используются слдеющие команды:
% - простая подпись: \sign[<сдвиг>]{<масштаб>}{<FIO>}
% - подпись с датой: \signat[<сдвиг>]{<масштаб>}{<FIO>}{<дата>}
% где
% - <сдвиг> --- необязательный сдвиг подписи по вертикали для правильного
%   расположения относительно строки
% - <масштаб> --- число, используемое для масштарибования изображения подписи
% - <FIO> --- имя файла с подписью, файл должен быть помещен в
%   img/signatures/FIO.png и иметь прозрачный фон
% - <дата> --- дата, которая будет выведена под подписью

% ## Утверждение задания руководителем и студентом
\authortaskapproval{}
\supervisortaskapproval{}

% ## Утверждение РСПЗ руководителем, студентом и консультантом
\authorrspzapproval{}
\supervisorrspzapproval{}
\consultantrspzapproval{}%

% ## Оценка руководителя за РСПЗ
\supervisorrspzgrade{}%

% ## Утверждение ПЗ руководителем, студентом и консультантом
\authorpzapproval{}%
\supervisorpzapproval{}
\consultantpzapproval{}%

% ## Оценка руководителя за ПЗ
\supervisorpzgrade{}%

%%% Local Variables:
%%% TeX-engine: xetex
%%% eval: (setq-local TeX-master (concat "../" (seq-find (-cut string-match ".*-1-task\.tex$" <>) (directory-files ".."))))
%%% End:
